\uuid{BlQS}
\titre{Logique}
\auteur{bodin ; croizat ; sacre}
\organisation{exo7}
\licence{CC BY-NC-SA 4.0}
\datecreate{2022-03-08}
\niveau{L1}
\module{Algèbre}
\chapitre{Logique, ensemble, raisonnement}
\sousChapitre{Logique}
\difficulte{2}
\theme{Logique}
\exo7qcmid{hilisit.logique.006}

\contenu{
\begin{question}
\qclassification{100.01}
On considère l'assertion "$\text{non}(\mathcal{P}) \text{ et } \mathcal{Q}$".
Quand est-ce que cette assertion est vraie ?
\begin{answers}
	\bad{Si $\mathcal{P}$ vraie et $\mathcal{Q}$ vraie.}
	\bad{Si $\mathcal{P}$ vraie et $\mathcal{Q}$ fausse.}
	\good{Si $\mathcal{P}$ fausse et $\mathcal{Q}$ vraie.}
	\bad{Si $\mathcal{P}$ fausse et $\mathcal{Q}$ fausse.}
\end{answers}
\begin{explanations}
Il faut  $\text{non}(\mathcal{P})$ vraie et $\mathcal{Q}$ vraie, c'est-à-dire
$\mathcal{P}$ fausse et $\mathcal{Q}$ vraie.
\end{explanations}
\end{question}
}
